% ============================================================================
% Supplementary material
% Uses the same preamble as the main paper (amsmath, amsthm, booktabs, graphicx,
% xcolor[table]) and the same labels (sec:theory, eq:floor, prop:pid, tab:main, ...).
% ============================================================================

\appendix
\renewcommand{\thesection}{S\arabic{section}}
\renewcommand{\thetable}{S\arabic{table}}
\renewcommand{\thefigure}{S\arabic{figure}}
\setcounter{table}{0}\setcounter{figure}{0}

\section*{Supplementary Material}

\noindent
The supplement contains the construction of every answer table (Sec.~\ref{sup:tables}),
the exact solver (Sec.~\ref{sup:milp}), full proofs (Sec.~\ref{sup:proofs}), details and
full results of the controlled study (Sec.~\ref{sup:controlled}), the VLM study
(Sec.~\ref{sup:vlm}) and the continual-learning study (Sec.~\ref{sup:cl}), and the analyses added for the revision:
uniqueness of exceptions (Sec.~\ref{sup:unique}), an accuracy-matched null for the VLM study (Sec.~\ref{sup:null}),
difficulty matching (Sec.~\ref{sup:difficulty}) and the size dependence of $s(T)$ (Sec.~\ref{sup:stability}). Code, pre-registrations and all result files are released with the paper.

% ----------------------------------------------------------------------------
\section{Answer tables}
\label{sup:tables}

Every table is computed from ground truth, never from a model. A question enters a table only
if it is well posed on every image of that table, so every cell has a defined answer.

\paragraph{CLEVR.}
Rows are ten yes/no questions of one CLEVR question family; columns are validation images. We
execute each question's functional program on each image's scene graph. Questions that refer
to ``the'' object are undefined on images without exactly one such object; we use set-lifted
semantics (``some'' object), which reproduces CLEVR's own answer on every question's original
image (60{,}835 of 60{,}835 yes/no questions). For the VLM study we keep only families whose
questions are well posed on every image under the natural reading (existence and count
comparisons; 24 units), so that a VLM is never asked an ill-posed question.

\paragraph{GQA.}
Rows are ten existence questions of one family (\texttt{exist}, \texttt{existC},
\texttt{existOr}, \texttt{existOrC}); columns are validation images containing at least one
object the questions mention. We execute each program on each scene graph with object
identifiers removed. A family is kept if execution reproduces GQA's own answers on at least
99\% of its questions, and questions about superordinate categories (``animal'', ``vehicle'')
are dropped because scene graphs label specific objects. This gives 24 units of 200 images.

\paragraph{CLOVE-scene.}
CLOVE~\cite{lei2023symbolic} splits GQA images into six scene tasks (shop and dining,
workplace, home or hotel, transportation, sport and leisure, outdoors) and releases the GQA
program of every question. We extend the executor with relations (\texttt{relate}) and keep
existence questions whose named objects are not superordinate categories and occur at least
50 times in GQA scene graphs. Family-level agreement with CLOVE's answers is 0.98--1.00. Units
take ten distinct programs of one family, 150 training images from the task's CLOVE train
split and 50 test images from its validation split, and require a yes-share between 0.10 and
0.90. Relation questions (``Is there bacon on the grill?'') are true on only about 1\% of
other images, even when sampling only images that contain all named objects, so every
relation unit fails the yes-share rule. The study uses 43 existence units, 6--8 per task.

\paragraph{POPE.}
Rows are COCO object categories and columns are the 500 POPE images; a cell is ``yes'' if the
object is in the image's COCO annotation, which reproduces all 9{,}000 POPE labels. For the VLM
study, eight units each take ten objects (shuffled) and 100 images.

\paragraph{Constructed tasks.}
Rows and columns are ten CIFAR-100 classes each. A sample takes one view of an image of class
$a$ and one view of a different image of class $b$, labelled $T(a,b)$. Tables are drawn from a
pool of 1{,}368 candidates spanning threshold tables, threshold tables with flipped cells, XOR
blocks and random tables.

% ----------------------------------------------------------------------------
\section{Exact computation of the additive rule}
\label{sup:milp}

With an error indicator $z_{ab}\in\{0,1\}$ per cell, a unit margin and a constant $M$,
\begin{align}
  \min_{u,v,z}\quad & \textstyle\sum_{(a,b)\in\Omega} w_{ab}\,z_{ab}\\
  \text{s.t.}\quad & u_a+v_b+M z_{ab}\ \ge\ 1 && \text{if } T(a,b)=1,\nonumber\\
  & u_a+v_b-M z_{ab}\ \le\ -1 && \text{if } T(a,b)=0,\nonumber\\
  & |u_a|,|v_b|\le U,\qquad z_{ab}\in\{0,1\}.\nonumber
\end{align}
Because the rule is unchanged by positive rescaling of $(u,v)$, the unit margin loses nothing
for $U$ and $M$ large enough; we use $U=50$, $M=201$. Columns with identical answers and
weights on every row can share one score, so we merge them and add their weights; a complete
table with ten rows then has at most $2^{10}$ distinct columns. We solve with HiGHS. When a
time limit is reached we keep the best rule and the solver's proven bound. In the VLM study,
95 of the 99 full tables are proven optimal; the other four (CLEVR count comparisons) are within
0.014 of the bound after one hour. Of the 99 confirmatory sub-tables, 98 are proven optimal. On
the full CLEVR set used for continual learning, 62\% of units are proven optimal and the mean
gap to the bound is 0.009. Finding the minimum number of cell flips that makes a table a Ferrers
relation is closely related to editing a bipartite graph into a chain graph, which is NP-hard in
general~\cite{yannakakis1981computing,natanzon2001complexity}; merging columns bounds the problem by the number of
rows, and every table here has ten. The solver returns one optimal rule; Sec.~\ref{sup:unique} examines the others.

\paragraph{Why not a trained additive model.}
Logistic regression on one-hot indices, any gradient-trained additive model, and EMAP optimise
a smooth surrogate (log-loss or squared error), not 0/1 accuracy. Their accuracy is a lower
bound on $C(T)$ and was 4--5 points below it on CLEVR and GQA.

% ----------------------------------------------------------------------------
\section{Proofs}
\label{sup:proofs}

\paragraph{Theorem~\ref{thm:ferrers}.}
(i)$\Rightarrow$(ii): if $u_a+v_b>0$, $u_{a'}+v_{b'}>0$, $u_a+v_{b'}\le 0$ and $u_{a'}+v_b\le 0$,
adding the first two and the last two gives $u_a+u_{a'}+v_b+v_{b'}$ both $>0$ and $\le 0$.
(ii)$\Rightarrow$(i): without an XOR square, the sets of ``yes'' columns of any two rows are
nested, so rows can be ordered by inclusion and columns by how many rows they answer ``yes'';
the table is then a staircase, which scores equal to these ranks reproduce exactly
(the Ferrers or biorder characterisation~\cite{ducamp1969composite,doignon1984realizable}).
(ii)$\Leftrightarrow$(iii): replacing an XOR square by its mirror image keeps every row and
column sum, so a table with an XOR square is not the only table with its sums. Conversely, by
Ryser's theorem~\cite{ryser1957combinatorial} any two $0/1$ tables with equal row and column
sums are connected by a sequence of such swaps; a table with no XOR square admits no swap, so
it is the only one. \qed

\paragraph{Proposition~\ref{prop:ceiling}.}
Given the indices, $x_1$ and $x_2$ are independent. Fix one draw $x_1^{(a)}$ per row and one
draw $x_2^{(b)}$ per column, and set $u_a=g_1(x_1^{(a)})$, $v_b=g_2(x_2^{(b)})$. This is an
additive rule, so its weighted accuracy is at most $C(T)$. For every cell, the pair
$(x_1^{(a)},x_2^{(b)})$ has the distribution of a test pair from that cell, so averaging over
draws gives exactly the expected accuracy of the readout. \qed

\paragraph{Proposition~\ref{prop:floor}.}
EMAP replaces the score matrix $\ell$ by $\hat\ell_{ab}=\bar\ell_{a\cdot}+\bar\ell_{\cdot b}-\bar\ell_{\cdot\cdot}$,
whose sign is the additive rule with $u_a=\bar\ell_{a\cdot}-\bar\ell_{\cdot\cdot}$ and
$v_b=\bar\ell_{\cdot b}$. Its error on the uniformly weighted complete table is at least $s(T)$.
With $e$ the error of $f$, $\Delta(f)=(1-e)-\mathrm{Acc}(\hat\ell)\ge(1-e)-(1-s(T))=s(T)-e$. \qed
The bound is tight for an exact rule-following model and vacuous when $e\ge s(T)$.

\paragraph{Proposition~\ref{prop:pid}.}
BROJA~\cite{bertschinger2014quantifying} defines unique information as
$\mathrm{UI}(y{:}a\setminus b)=\min_{q\in\Delta_p} I_q(y;a\mid b)$ over the set $\Delta_p$ of joint
distributions with the same $p(a,y)$ and $p(b,y)$; redundancy and synergy follow from UI, the
pairwise informations $I(y;a)$, $I(y;b)$ (functions of $p(a,y)$, $p(b,y)$) and the total
$I(y;a,b)$. For deterministic answers $I(y;a,b)=H(y)$, a function of $p(y)$. Every term is
therefore a function of $p(a,y)$ and $p(b,y)$ alone. \qed

\paragraph{Example on partial tables.}
On partial tables PID cannot even detect whether exceptions exist. The two $3\times3$ tables
below (``$\cdot$'' = cell not in $\Omega$, uniform weights) have identical row and column
answer counts and therefore identical BROJA atoms (synergy 0.667, redundancy 0.126, unique
0.000 and 0.126 bits), yet the first is exactly additive and the second contains an XOR square.
\begin{center}\small
\begin{tabular}{ccc}0&1&0\\0&$\cdot$&1\\$\cdot$&0&$\cdot$\end{tabular}
\qquad\qquad
\begin{tabular}{ccc}0&1&0\\$\cdot$&0&1\\0&$\cdot$&$\cdot$\end{tabular}
\end{center}

\paragraph{BROJA solver.}
For the 10$\times$100 VLM tables we compute BROJA with the convex solver of
Liang~\etal~\cite{liang2023quantifying} (SCS). On six 10$\times$20 checks it matches the exact
\texttt{dit} implementation to within 0.0003 bits.

% ----------------------------------------------------------------------------
\section{Controlled study}
\label{sup:controlled}

From the pool of 1{,}368 tables we select 63 with BROJA synergy between 0.40 and 0.75 bits
such that synergy and XOR density are nearly uncorrelated ($+0.17$). Each modality is a frozen
image feature; the full model fuses two small encoders through an MLP head, and the linear
readout uses the same features. Three seeds per table. The gap is the full model's accuracy
minus the readout's, as a share of the full model's accuracy above the majority baseline.

\begin{table*}[h]
\centering\small
\caption{Controlled tables. Partial rank correlations with the readout gap, each controlling
for the other predictor; 95\% bootstrap intervals over tables.}
\label{tab:sup_controlled}
\begin{tabular}{lcc}
\toprule
Test & Structure & PID synergy \\
\midrule
63 tables, synergy and structure decorrelated & $+0.91$ $[0.83, 0.94]$ & $+0.22$ $[-0.06, 0.45]$ \\
62 fresh tables (replication) & $+0.86$ $[0.74, 0.93]$ & $+0.11$ $[-0.14, 0.36]$ \\
Repair after drift (125 tables) & $+0.86$ $[0.80, 0.90]$ & $-0.10$ $[-0.28, 0.09]$ \\
12 families with identical PID & $+0.78$ (within family) & constant \\
\bottomrule
\end{tabular}
\end{table*}

\noindent
The same holds with $I_{\mathrm{ccs}}$ ($+0.85$ vs.\ $+0.25$) and MMI ($+0.91$ vs.\ $+0.22$)
in place of BROJA. All 16 tables without exceptions have a gap of at most 0.09 (median 0.06)
while carrying up to 0.67 bits of synergy. One pre-registered comparison went the other way:
for accuracy on held-out cell combinations, synergy predicted slightly better than structure
($+0.36$ vs.\ $+0.31$).

% ----------------------------------------------------------------------------
\section{VLM study}
\label{sup:vlm}

\paragraph{Models and inference.}
Qwen2.5-VL-7B-Instruct, InternVL3.5-8B (HF port), LLaVA-OneVision-1.5-8B-Instruct (publisher
code at a pinned revision, run with \texttt{transformers} 4.53, the version it targets) and
Idefics3-8B-Llama3. Images are scaled to a longest side of 448 pixels; InternVL3.5 and
Idefics3 use a single image tile instead of their default multi-tile cropping. The prompt is
``\textit{$\langle$question$\rangle$} Answer with yes or no.'' and the score is the log-odds of
the first token of ``Yes/yes'' versus ``No/no'' at the next position. No model is fine-tuned.

\paragraph{EMAP.}
For each complete 10$\times$100 table we compute $\hat\ell$ from the model's log-odds on all
1{,}000 cells and report the accuracy of $\ell$ minus that of $\hat\ell$.

\paragraph{Study A: full tables.}
Within each benchmark, the exception share and BROJA synergy are nearly collinear (Spearman
$+0.95$ CLEVR, $+0.90$ CLOVE, $+0.88$ GQA, $+0.48$ POPE), so the two cannot be separated.

\begin{table*}[h]
\centering\small\setlength{\tabcolsep}{3pt}
\caption{Study A. Within-benchmark Spearman correlation of the EMAP drop with the exception
share (synergy in parentheses).}
\begin{tabular}{lcccc}
\toprule
Model & CLEVR & GQA & CLOVE & POPE \\
\midrule
Qwen2.5-VL-7B   & $+0.37$ $(+0.34)$ & $+0.07$ $(-0.08)$ & $+0.61$ $(+0.48)$ & $+0.44$ $(+0.78)$ \\
InternVL3.5-8B  & $+0.91$ $(+0.90)$ & $+0.11$ $(+0.11)$ & $+0.76$ $(+0.70)$ & $+0.49$ $(+0.51)$ \\
LLaVA-OV-1.5-8B & $+0.97$ $(+0.96)$ & $+0.09$ $(+0.17)$ & $+0.58$ $(+0.47)$ & $+0.48$ $(+0.85)$ \\
\bottomrule
\end{tabular}
\end{table*}

\paragraph{Selecting decorrelated sub-tables.}
From each table we draw 27 sub-tables of 50 images (2{,}673 candidates). Image weights are
tilted towards or away from columns rich in exceptions or in ``yes'' answers,
$w_b\propto\exp(\alpha z_b^{\mathrm{exc}}+\beta z_b^{\mathrm{yes}})$ with
$\alpha,\beta\in\{-2,0,2\}$ and three draws each, so that structure and synergy vary partly
independently. Each candidate gets its exact rule (MILP, 300 s; 97.5\% proven optimal), BROJA
and MMI. One candidate per table is then chosen by coordinate descent (30 sweeps) to minimise,
within each benchmark, $|\rho(s,\mathrm{BROJA})|$ (study B) or
$|\rho(s,\mathrm{BROJA})|+|\rho(s,\mathrm{yes})|$ (study C), while keeping the spread of $s$ at
least half that of the candidate pool. The selection uses answer tables only. Achieved
correlations: study B, $\rho(s,\mathrm{BROJA})=0.00$ in every benchmark; study C,
$\rho(s,\mathrm{BROJA})\le 0.02$ (GQA $0.20$) and $\rho(s,\mathrm{yes})\le 0.00$.

\begin{table}[t]
\centering
\caption{Partial rank correlations of EMAP drop. Study B controls for the other measure and accuracy. Study C additionally controls for yes-share. Values are correlations with 95\% bootstrap intervals over tables.}
\label{tab:partial_corr}
\resizebox{\columnwidth}{!}{%
\begin{tabular}{llcc}
\toprule
Study & Model & Structure & Synergy \\
\midrule
B & Qwen2.5-VL-7B & $+0.33$ $[0.10,0.50]$ & $+0.19$ \\
B & InternVL3.5-8B & $+0.37$ $[0.17,0.54]$ & $+0.44$ \\
B & LLaVA-OV-1.5-8B & $+0.32$ $[0.09,0.50]$ & $+0.43$ \\
\midrule
C & Idefics3-8B & $+0.42$ $[0.21,0.56]$ & $+0.07$ $[-0.17,0.29]$ \\
C & InternVL3.5-8B & $+0.42$ $[0.20,0.58]$ & $+0.05$ $[-0.16,0.26]$ \\
C & LLaVA-OV-1.5-8B & $+0.43$ $[0.23,0.58]$ & $+0.24$ $[0.02,0.42]$ \\
C & Qwen2.5-VL-7B & $+0.26$ $[0.06,0.44]$ & $+0.03$ $[-0.18,0.25]$ \\
\bottomrule
\end{tabular}%
}
\end{table}


\noindent
In study B, synergy kept a signal of its own. A post-hoc check found that, once structure was
fixed, synergy followed the yes-share ($+0.57$ to $+0.83$); adding the yes-share as a control
removed synergy's signal and strengthened structure's ($+0.37$ to $+0.48$). Study C was
pre-registered to test this on a new model and a fresh selection. Its secondary linear
regression (standardised, with benchmark effects) failed for all four models: on the raw
scale, structure and synergy each add only 1--4\% of explained variance beyond benchmark and
answer balance (Idefics3: $+0.014$ and $+0.021$).

\paragraph{Floor and cells.}
As a sanity check of the pipeline, $\Delta\ge s-e$ held for all $4\times99$ model--table pairs, using the solver's
proven lower bound on $s$ where the rule was not proven optimal. This is guaranteed by Prop.~\ref{prop:floor} and is
not evidence for the measure.

\begin{table}[t]
\centering
\caption{Cell-level AUC for predicting the answers EMAP flips (full tables, three models). Binary exception labels use
the solver's rule; Sec.~\ref{sup:unique} reports stable labels on the Study C sub-tables, where the binary label
($0.60$--$0.66$) does not beat confidence ($0.63$--$0.73$) but graded misfit does ($0.78$--$0.88$).}
\label{tab:cell_auc}
\resizebox{\columnwidth}{!}{%
\begin{tabular}{lccc}
\toprule
Model & Exception & Graded misfit & Low confidence \\
\midrule
Qwen2.5-VL-7B    & 0.68 & 0.87 & 0.74 \\
InternVL3.5-8B    & 0.62 & 0.78 & 0.67 \\
LLaVA-OV-1.5-8B  & 0.70 & 0.88 & 0.64 \\
\bottomrule
\end{tabular}%
}
\end{table}


% ----------------------------------------------------------------------------
\section{Continual-learning study}
\label{sup:cl}

\paragraph{Network and training.}
Each input is a frozen feature vector: CIFAR image features for constructed tasks, symbolic
scene features and CLIP question features for CLEVR, and CLIP image and question features for
GQA and CLOVE. Each modality has a trainable MLP encoder (two hidden layers of 256 units,
128-dimensional output); the fusion head is an MLP on the concatenated codes. Adam, learning rate
$10^{-3}$, weight decay $10^{-5}$, batch 256, 50{,}000 training draws per task. The first task
is trained for 30 epochs and later tasks for 10. Buffers hold 1{,}000 samples (split equally
over past tasks in the six-task setting), with 32 replayed samples per step; DER++ uses
$\alpha=0.5$.

\paragraph{Matched plasticity.}
Strengths are chosen per run from EWC $\{1,10^2,10^4\}$, SI $\{0.1,1,10,100\}$, LwF and
feature distillation $\{0.3,3\}$, and iCaRL distillation $\{0.1,1\}$: among settings whose
accuracy on the new task is within 0.02 of finetuning, the one that best keeps the old task.

\paragraph{Benchmarks.}
Constructed: 93 tables with exceptions, each followed by a balanced table on other classes.
CLEVR: 40 units, each followed by a unit of another family. GQA: 24 units with three seeds,
each followed by a unit of another family. CLOVE: six scene tasks in CLOVE's six released
orders, two seeds (12 runs).
\begin{figure}[t]\centering
\includegraphics[width=\linewidth]{../../figures/revision/fig_dbwt_cb.pdf}
\caption{Backward-transfer gap between exceptions and rule cells with the same answer, with 95\% intervals; filled:
interval below zero. Gaps are only meaningful for methods that keep the rule: finetuning and weak regularisers show
positive gaps because exceptions start near chance.}
\label{fig:dbwt}\end{figure}
\paragraph{Same-answer gaps.}
With $\mathrm{BWT}_{g,y}$ the backward transfer on cells of type $g\in\{E,R\}$ (exception or
rule) and answer $y$, and $\pi_y$ the share of exceptions with answer $y$,
\begin{equation}
  \Delta\mathrm{B}=\sum_{y}\pi_y\left[\mathrm{BWT}_{E,y}-\mathrm{BWT}_{R,y}\right],
\end{equation}
and $\Delta\mathrm{F}$ likewise from forward transfer against a model trained on the new task
alone. Answers with no cells of one type are dropped and $\pi$ is renormalised. Intervals are
bootstraps over units.

\paragraph{Why matching answers matters.}
Exceptions are 92\% ``yes'' on GQA and 89\% on CLOVE, against 6\% and 8\% of rule cells.
Without matching, every protective method appeared to lose exceptions faster on all
benchmarks; with matching, the backward effect remains strong only on CLEVR (Tab.~\ref{tab:main}).
Our first, unmatched analyses were pre-registered and passed; we report the matched analysis
as the main result because the unmatched one partly measures class imbalance.
\begin{figure}[t]\centering
\includegraphics[width=\linewidth]{../../figures/revision/fig_buffers_cb.pdf}
\caption{\textbf{What a replay buffer should hold.} Four DER++ buffers of 1{,}000 samples. Arrow and number: exception
BWT gained by filling half the buffer with exception cells at fixed answer balance (pre-registered). The gain is
7--9 points except on GQA, where answer balance alone explains it, and costs 1--3 points of overall BWT.}
\label{fig:buffers}\end{figure}
\paragraph{Buffer control.}
DER++ with four buffers of 1{,}000 samples (pre-registered).

\begin{table}[t]
\centering
\caption{Backward transfer on exceptions / overall for four DER++ buffers.}
\label{tab:der_buffers}
\resizebox{\columnwidth}{!}{%
\begin{tabular}{lcccc}
\toprule
Buffer & Constructed & CLEVR & GQA & CLOVE \\
\midrule
Random
& $-.034$ / $-.021$ & $-.085$ / $.000$ & $-.040$ / $.004$ & $-.153$ / $-.007$ \\
Class balanced
& $-.025$ / $-.022$ & $-.059$ / $-.001$ & $-.009$ / $.001$ & $-.007$ / $-.018$ \\
Exception focused
& $.058$ / $-.032$ & $.001$ / $-.006$ & $-.007$ / $.001$ & $.054$ / $-.026$ \\
Exception focused, balanced
& $.053$ / $-.048$ & $.008$ / $-.008$ & $-.014$ / $.001$ & $.077$ / $-.035$ \\
\bottomrule
\end{tabular}%
}
\end{table}

\noindent
Balanced rule+exceptions versus balanced random, exception BWT: $+0.078$ $[0.069, 0.088]$,
$+0.068$ $[0.054, 0.081]$, $-0.005$ $[-0.016, 0.005]$, $+0.085$ $[0.058, 0.112]$ (pass on 3 of
4). The overall cost bound ($>-0.02$) failed on constructed tasks ($-0.027$) and CLOVE
($-0.028$).

\paragraph{Method guide (negative result).}
We tested whether data descriptors (exception share, answer skew, yes-share, exception
accuracy after learning) predict which method keeps the most exceptions on a held-out
benchmark. The descriptors beat a fixed method ranking on 3 of 4 held-out benchmarks but did
not beat the standard ranking by overall backward transfer (1 of 4).
\begin{figure}[t]\centering
\includegraphics[width=\linewidth]{../../figures/revision/fig_bwt_exc.pdf}
\caption{\textbf{Standard backward transfer picks the wrong winner} (CLEVR). DER++ has the best overall BWT but the
lowest exception accuracy among replay methods, below chance; DER-expansion keeps both. Open grey circle: the
DER++ buffer with exception cells, a diagnostic probe.}
\label{fig:bwt_exc}\end{figure}
\paragraph{VLM continual learning.}
Qwen2.5-VL-3B-Instruct with LoRA ($r=16$, $\alpha=32$) on all language-model projections;
images 336$\times$336; 300 steps of 8 examples per task, learning rate $10^{-4}$; memory of
1{,}000 examples with 4 replayed per step. On CLOVE, the pre-registered tests gave: exceptions
forgotten faster for LwF and ER (pass, 2 of 3 methods, unmatched); DER++ forward gap $-0.029$
(fail, threshold $-0.03$); exception buffer (fail, DER++ already kept exceptions).
\begin{figure}[t]\centering
\includegraphics[width=\linewidth]{../../figures/revision/fig_decor.pdf}
\caption{\textbf{Separating structure from synergy.} Exception share against BROJA synergy, as ranks within each
benchmark. \textbf{(a)} On the full tables the two are nearly collinear. \textbf{(b)} The confirmatory sub-tables,
chosen from answer tables only, make them uncorrelated, and also uncorrelated with the yes share ($|\rho|<0.01$).}
\label{fig:decor}\end{figure}

\paragraph{Nonce-word CLEVR.}
On plain CLEVR, Qwen2.5-VL-3B already answers 92--98\% of cells zero-shot, rule and exception
cells alike, so a task sequence teaches it nothing. We therefore rename the 15 attribute
concepts (eight colours, three shapes, two materials, two sizes, with synonyms and plurals) in
each task with that task's own made-up words; answers and exception labels are unchanged. A
pilot set the budget (300 steps). The pre-registered learning check (task-only accuracy at
least 0.10 above zero-shot) failed: 0.651 versus 0.587. The reason is that rule cells are
learned (0.64 to 0.86) while exceptions get worse (0.53 to 0.44), and the check was measured on
test sets that are half exceptions. We report the experiment as descriptive.
\begin{figure}[t]\centering
\includegraphics[width=\linewidth]{../../figures/revision/fig_vlm_trade.pdf}
\caption{\textbf{VLMs trade exceptions for the rule.} \textbf{(a)} Qwen2.5-VL-3B with LoRA on six CLOVE-scene tasks:
rule-cell accuracy rises slightly while exception accuracy falls from 82\% zero-shot to 66--72\% after every
method (dashed: DER++ with exception cells in its buffer). \textbf{(b)} With made-up attribute words, LoRA learns the
rule but exceptions get worse, below chance (descriptive; the pre-registered learning check failed).}
\label{fig:vlm_trade}\end{figure}
\paragraph{Pretraining as the first task (post hoc).}
Measuring change from the zero-shot model, with answers matched: fine-tuning on the related
CLOVE sequence lowers exceptions more than rule cells ($-0.07$); fine-tuning on an unrelated
OCR task does not ($+0.04$, mixed by answer).

% ----------------------------------------------------------------------------

% ----------------------------------------------------------------------------
\section{Uniqueness of exceptions}
\label{sup:unique}

The exception share is unique; the exception cells are not. We re-solve every table with random tie-breaking: the
objective adds $\varepsilon r_{ab} z_{ab}$ with $r_{ab}\sim U(0,1)$ per merged cell and $\varepsilon$ small enough that
the total perturbation is below one error, so every solution found has the minimum number of errors (20 solves per
table; 10 for the $10\times300$ CLEVR test tables, keeping only solutions with the smallest error count). A cell is a
\emph{stable} exception (rule cell) if it is an exception (rule cell) in every accepted solution and \emph{ambiguous}
otherwise. Because tie-breaking samples the optimal rules, the ambiguous share is a lower bound.

\begin{table}[h]
\centering\small\setlength{\tabcolsep}{3pt}
\caption{Ambiguity of exception cells. Saved: labels from the solver's single rule, used in the main analyses.}
\label{tab:unique}
\begin{tabular}{lccccc}
\toprule
Tables & $n$ & ambiguous & stable exc. & $s(T)$ & saved exc.\ that are stable \\
\midrule
VLM, Study C        & 99 & 0.054 & 0.042 & 0.067 & -- \\
CL, constructed     & 94 & 0.058 & 0.037 & 0.058 & 0.66 \\
CL, CLEVR test      & 64 & 0.032 & 0.074 & 0.090 & 0.75 \\
CL, GQA test        & 24 & 0.054 & 0.045 & 0.070 & 0.64 \\
CL, CLOVE test      & 43 & 0.048 & 0.036 & 0.060 & -- \\
\bottomrule
\end{tabular}
\end{table}

\noindent
Saved rule cells are stable rule cells in 97\% of cases. One table is an outlier: the constructed tasks' second table,
a random balanced $10\times10$ table, has every cell ambiguous across 19 optimal rules (29 errors each), so forward
gaps on constructed tasks depend on an arbitrary choice of rule and are omitted from Tab.~\ref{tab:main}.

\begin{table}[h]
\centering\small\setlength{\tabcolsep}{3pt}
\caption{Same-answer gaps (points) with the saved labels and with stable labels only (ambiguous cells dropped);
pre-registered robustness checks.}
\label{tab:unique_cl}
\begin{tabular}{llcc}
\toprule
 & Method & saved & stable \\
\midrule
CLEVR, backward & ER     & $-8.7$ $[-11.5,-5.8]$  & $-9.3$ $[-12.6,-5.9]$ \\
                & A-GEM  & $-11.0$ $[-14.2,-7.2]$ & $-12.0$ $[-15.7,-8.4]$ \\
                & GEM    & $-10.0$ $[-13.7,-5.6]$ & $-10.7$ $[-15.2,-5.8]$ \\
                & DER++  & $-9.2$ $[-10.6,-7.5]$  & $-10.2$ $[-11.8,-8.3]$ \\
CLEVR, forward  & DER++  & $-25.6$ $[-29.5,-19.8]$ & $-27.8$ $[-31.8,-22.1]$ \\
                & Feat.\ dist. & $-13.2$ $[-16.1,-9.8]$ & $-14.8$ $[-17.7,-11.0]$ \\
GQA, forward    & DER++  & $-7.1$ $[-12.3,-1.5]$  & $-6.0$ $[-12.6,1.2]$ \\
                & Feat.\ dist. & $-7.5$ $[-13.3,-1.9]$ & $-6.8$ $[-13.8,0.8]$ \\
\bottomrule
\end{tabular}
\end{table}

\noindent
Pre-registered verdicts: the CLEVR backward result holds for 4 of 4 replay methods (pass); the DER++ forward result
holds on 1 of 3 benchmarks (fail: CLEVR holds, the GQA interval reaches zero, and constructed tasks have no stable
labels). On the Study C sub-tables, stable exception labels predict the cells EMAP flips with AUC $0.60$--$0.66$, below
or near model confidence ($0.63$--$0.73$), while graded misfit, which needs no choice of rule, reaches $0.78$--$0.88$.

% ----------------------------------------------------------------------------
\section{Accuracy-matched null for the VLM study}
\label{sup:null}

For each model and Study C sub-table, the null scorer has logits $(2T_{ab}-1)\mu_y+\varepsilon_{ab}$,
$\varepsilon_{ab}\sim\mathcal N(0,1)$, with $\mu_y=\Phi^{-1}(\mathrm{acc}_y)$ matching the model's accuracy on cells with
answer $y$ (mean absolute accuracy difference $0.002$); its EMAP drop $\Delta_0$ is averaged over 20 draws.

\begin{table}[h]
\centering\small\setlength{\tabcolsep}{2.5pt}
\caption{Within-benchmark partial rank correlations (controls: synergy, accuracy, yes-share unless stated); 95\%
bootstrap intervals over tables. Adjusted: $\Delta$ with $\Delta_0$ as an extra control. Yes: partial correlation of
$\Delta$ with the yes-share (controls: $s$, synergy, accuracy).}
\label{tab:sup_null}
\begin{tabular}{lcccccc}
\toprule
Model & $\bar\Delta$ / $\bar\Delta_0$ & $\rho(\Delta,s)$ & $\rho(\Delta_0,s)$ & $\rho(\Delta-\Delta_0,s)$ & adjusted & yes \\
\midrule
Idefics3     & .064 / .060 & $+.42$ & $+.37$ & $+.17$ $[-.07,.38]$ & $+.25$ $[.01,.44]$  & $+.23$ $[.03,.41]$ \\
InternVL3.5  & .086 / .044 & $+.42$ & $+.50$ & $+.26$ $[.05,.43]$  & $+.20$ $[-.03,.41]$ & $+.34$ $[.14,.50]$ \\
LLaVA-OV-1.5 & .081 / .075 & $+.43$ & $+.40$ & $+.12$ $[-.11,.32]$ & $+.23$ $[-.01,.41]$ & $+.30$ $[.10,.49]$ \\
Qwen2.5-VL   & .079 / .070 & $+.26$ & $+.30$ & $+.05$ $[-.17,.25]$ & $+.08$ $[-.11,.29]$ & $+.32$ $[.09,.49]$ \\
\bottomrule
\end{tabular}
\end{table}

\noindent
Pre-registered primary test (Idefics3, $\rho(\Delta-\Delta_0,s)$ lower bound above zero): fail. The share of the
observed correlation that survives subtracting the null is 42\% (Idefics3), 63\% (InternVL3.5), 28\% (LLaVA-OV-1.5) and
19\% (Qwen2.5-VL). The slack $\Delta-(s-e)$ correlates with $s$ only for InternVL3.5 ($+0.28$ $[0.06,0.45]$).

% ----------------------------------------------------------------------------
\section{Difficulty-matched continual learning}
\label{sup:difficulty}

We re-ran finetuning, LwF, feature distillation, ER, DER++ and DER-expansion on the constructed tasks, CLEVR and GQA
with the same units, seeds and first-task training, saving each test cell's signed margin for its true answer: the
first-task model's before the second task (forgetting) and a second-task-only model's (learning). Strata are unit
$\times$ answer $\times$ margin decile (deciles within unit and answer); each stratum with at least three exceptions and
three rule cells contributes its exceptions-minus-rule difference, weighted by its number of exceptions. Matched rule
cells exist for 99--100\% of exceptions. We also compare exceptions with the hardest quarter of rule cells of the same
answer.

\begin{table}[h]
\centering\small\setlength{\tabcolsep}{3pt}
\caption{Constructed tasks, backward gaps (points). Forward gaps are not interpretable here (Sec.~\ref{sup:unique}).}
\label{tab:sup_diff_tables}
\begin{tabular}{lccc}
\toprule
Method & same answer & + difficulty & vs hardest quarter \\
\midrule
Finetune      & $+9.7$  & $+2.0$ $[0.0,3.8]$    & $-22.5$ \\
LwF           & $-0.5$  & $-2.3$ $[-4.3,-0.5]$  & $-8.3$ \\
Feature dist. & $+11.0$ & $+3.7$ $[1.0,6.4]$    & $-20.9$ \\
ER            & $-0.6$  & $-2.8$ $[-4.1,-1.4]$  & $-12.7$ \\
DER++         & $-1.0$  & $-1.9$ $[-2.9,-0.9]$  & $-3.6$ \\
DER-expansion & $+0.5$  & $-0.5$ $[-1.6,0.6]$   & $-4.6$ \\
\bottomrule
\end{tabular}
\end{table}

\noindent
CLEVR is in Tab.~\ref{tab:difficulty}; on CLEVR exceptions also fare worse than the hardest quarter of rule cells for
every method (backward $-6.7$ to $-20.6$, forward $-21.7$ to $-42.6$ points). On GQA (Tab.~\ref{tab:sup_diff_gqa}) only 18--19\% of exceptions have matched rule cells, so its
difficulty-matched estimates rest on a small, selected part of the exceptions.

\begin{table}[h]
\centering\small\setlength{\tabcolsep}{3pt}
\caption{GQA, difficulty-matched gaps (points; coverage 0.18--0.19).}
\label{tab:sup_diff_gqa}
\begin{tabular}{lcccc}
\toprule
 & \multicolumn{2}{c}{Backward} & \multicolumn{2}{c}{Forward} \\
Method & same answer & + difficulty & same answer & + difficulty \\
\midrule
Finetune      & $-0.6$ & $-11.4$ $[-25.1,1.2]$ & $-3.9$ & $-0.9$ $[-7.7,5.4]$ \\
LwF           & $+4.9$ & $-6.2$ $[-16.4,4.4]$  & $-2.5$ & $-9.8$ $[-17.1,-2.8]$ \\
Feature dist. & $+0.2$ & $-16.0$ $[-30.7,-0.7]$ & $-5.5$ & $-3.5$ $[-10.3,2.0]$ \\
ER            & $-1.5$ & $-4.0$ $[-10.0,2.8]$  & $-2.4$ & $-1.2$ $[-9.8,6.8]$ \\
DER++         & $-2.8$ & $-2.5$ $[-6.3,0.9]$   & $-3.3$ & $-9.5$ $[-18.6,-0.0]$ \\
DER-expansion & $-3.0$ & $-4.6$ $[-8.8,-0.9]$  & $-1.4$ & $-3.0$ $[-11.3,3.7]$ \\
\bottomrule
\end{tabular}
\end{table}

\noindent
Pre-registered verdicts: D2 (CLEVR backward, ER and DER++ $\le-3$ points with interval below zero) pass; D1 (DER++
forward $\le-3$ points with interval below zero on at least two of constructed tasks, CLEVR and GQA) pass, 3 of 3.
D1 should be read with two caveats found by the other revision analyses: the constructed tasks' forward labels are not
unique (Sec.~\ref{sup:unique}), and on GQA the interval only just excludes zero at 18\% coverage. Only CLEVR is
unaffected by both.

% ----------------------------------------------------------------------------
\section{Size dependence of the exception share}
\label{sup:stability}

For each CLEVR continual-learning unit ($10\times1{,}000$) and each GQA and CLOVE unit ($10\times200$) we draw column
subsets of $n$ images (five draws each) with all ten questions or five random questions and compute $s$ exactly.

\begin{table}[h]
\centering\small\setlength{\tabcolsep}{3pt}
\caption{Mean exception share by number of images $n$ (10 questions / 5 questions).}
\label{tab:sup_stability}
\begin{tabular}{lcccccc}
\toprule
 & 25 & 50 & 100 & 200 & 400 & 1000 \\
\midrule
CLEVR & .079/.059 & .086/.063 & .090/.066 & .093/.070 & .096/.069 & .099/.072 \\
GQA   & .065/.048 & .069/.052 & .070/.052 & .071/.048 & -- & -- \\
CLOVE & .059/.040 & .060/.040 & .061/.046 & .062/.037 & -- & -- \\
\bottomrule
\end{tabular}
\end{table}

\noindent
The share grows slowly with the number of images and is lower with fewer questions, so absolute values should be
compared at a fixed table shape. Relative values are stable: 84--97\% of the variance in $s$ is between tables rather
than between draws of the same table (at $n=50$ and $100$), and the ranking of tables at $n=50$ matches the full-size
ranking (Spearman $0.97$--$0.99$).

% ----------------------------------------------------------------------------
\section{Pre-registration record}
\label{sup:prereg}

Before each study we wrote its hypotheses, thresholds and analysis code and recorded SHA-256 hashes of the files.
[[Anonymised registration link (e.g.\ OSF, with timestamps) to be added.]] The revision analyses (null scorer,
uniqueness, difficulty matching, size dependence) were pre-registered together (hash \texttt{73d425b2}): the null test
failed, the uniqueness tests passed for forgetting and failed for learning, and the difficulty tests passed for
forgetting on CLEVR and, formally, for learning (robustly only on CLEVR; Sec.~\ref{sup:difficulty}). Deviations: the GPU driver of our machine changed, so GPU work used the
same PyTorch version (2.14.0) built for CUDA 12.6; the GQA difficulty runs crashed twice on library builds and were
re-run with matching builds and unchanged code.
